% Generated draft: requires mathematical and source review.
\documentclass[11pt]{article}
\usepackage[margin=25mm]{geometry}
\usepackage{fontspec}
\setmainfont{Latin Modern Roman}
\usepackage{amsmath,amssymb,mathtools,unicode-math}
\setmathfont{Latin Modern Math}
\usepackage{enumitem,needspace,xurl,hyperref}
\setlength{\parindent}{0pt}
\setlength{\parskip}{5pt}
\setlength{\emergencystretch}{3em}
\newcommand{\N}{\mathbb{N}}
\newcommand{\Z}{\mathbb{Z}}
\newcommand{\Q}{\mathbb{Q}}
\newcommand{\R}{\mathbb{R}}
\newcommand{\C}{\mathbb{C}}
\newcommand{\F}{\mathbb{F}}
\newcommand{\A}{\mathbb{A}}
\newcommand{\PP}{\mathbb{P}}
\newcommand{\E}{\mathbb{E}}
\newcommand{\eps}{\varepsilon}
\newcommand{\dd}{\,\mathrm{d}}
\begin{document}
\section*{UnsolvedMath ID: 701617}
\textbf{Source problem code:} catalogue entry
\subsection{Definitions and mathematical statement}
A topological group is called \emph{tdlc} if it is totally disconnected and locally compact, and \emph{tdlcsc} if it is additionally second countable. We use the usual Hausdorff convention for locally compact groups; the supplied record does not explicitly state this convention. For a topological group \(H\), define its \emph{discrete residual} by
\[
\operatorname{Res}(H)=\bigcap\{N\leq H:N\text{ is open and normal in }H\}.
\]
A profinite group is an inverse limit of finite discrete groups. A group is compactly generated if it is generated by a compact subset.

Let \(\mathscr{E}\) be the smallest class of tdlcsc groups containing every second countable profinite group and every countable discrete group, and closed under closed subgroups, Hausdorff quotients, directed unions of open subgroups, and group extensions. Here closure under extensions means that a tdlcsc group \(G\) belongs to \(\mathscr{E}\) whenever it has a closed normal subgroup \(N\) with \(N,G/N\in\mathscr{E}\). For directed unions, the constituent open subgroups are directed by inclusion, and the resulting group must remain tdlcsc.

The supplied source describes an ordinal-valued decomposition rank \(\xi\) on \(\mathscr{E}\), with
\[
\xi(\{1\})=1.
\]
For a nontrivial \(G\in\mathscr{E}\), if
\[
G=\bigcup_{i\geq 1}O_i,\qquad O_i\leq O_{i+1},
\]
where each \(O_i\) is a compactly generated open subgroup of \(G\), the stated recursive rule is
\[
\xi(G)=\left(\sup_{i\geq 1}\xi\bigl(\operatorname{Res}(O_i)\bigr)\right)+1,
\]
where \(+1\) denotes ordinal successor. The source asserts that this rank is well behaved; its well-definedness and independence of the exhaustion are taken as supplied background. The restriction to nontrivial \(G\) makes explicit that the base case takes precedence: the source writes the recursive rule without this exception, although applying it to \(G=\{1\}\) would give \(2\) instead of \(1\).

\textbf{Problem.} Determine the ordinal
\[
S=\sup\{\xi(G):G\in\mathscr{E}\}.
\]
In particular, is \(S\) countable, equivalently, is \(S<\omega_1\), where \(\omega_1\) is the first uncountable ordinal?
\subsection{Short English statement}
Determine the supremum of the decomposition ranks of all elementary totally disconnected, locally compact, second countable groups. In particular, decide whether that supremum is a countable ordinal. The trivial-group base case is given precedence over the recursive rule, whose supplied wording omits this exception.
\subsection{Sources}
\begin{itemize}
\item Local UnsolvedMath registry, supplied complete record with permanent ID \(701617\), exact source code \texttt{catalogue entry}, titled \emph{Kourovka Notebook Problem 19.69}; upstream revision \texttt{d371ee917ac9647fbb06cd6fc529deec6f3dda82}. The supplied record labels the problem \emph{open}; this is a registry label, not an independently verified status assessment.
\item \emph{The Kourovka Notebook (Unsolved Problems in Group Theory)}, Problem 19.69. The supplied background cites the 19th issue (2018) and the 21st issue edition (2026), identified as \texttt{arXiv:1401.0300v48}. The collection description names E. I. Khukhro and V. D. Mazurov as editors. These bibliographic details are supplied attribution and have not been independently checked.
\item Attribution is unresolved: the record's proposer field and background name C. Reid, with proposed year 2018, whereas the mathematical statement begins with the attribution ``(P. Wesolek).'' Both are retained without assigning their respective roles.
\end{itemize}
\end{document}
